% Source: https://artofproblemsolving.com/community/c3h1669230_200th_post_problem
% Stable identity: aops-topic-1669230-q01
If $n=p_1^{e_1}p_2^{e_2}\cdots$, then $n$ has $(e_1+1)(e_2+1)\cdots$ divisors. For the smallest $n$ with a given list of exponents, the largest exponents go on the smallest primes. So we must write $200=2^3\cdot5^2$ as a product of factors $f_1\ge f_2\ge\dots\ge2$ and minimise $2^{f_1-1}3^{f_2-1}5^{f_3-1}\cdots$. Since $200$ has five prime factors, at most five factors appear. The main candidates are
\[
\begin{array}{ll}
5\cdot5\cdot2\cdot2\cdot2: & 2^4\cdot3^4\cdot5\cdot7\cdot11=498960\\
5\cdot5\cdot4\cdot2: & 2^4\cdot3^4\cdot5^3\cdot7=1134000\\
10\cdot5\cdot2\cdot2: & 2^9\cdot3^4\cdot5\cdot7=1451520\\
10\cdot5\cdot4: & 2^9\cdot3^4\cdot5^3=5184000\\
8\cdot5\cdot5: & 2^7\cdot3^4\cdot5^4=6480000
\end{array}
\]
and every other factorisation ($10\cdot10\cdot2$, $20\cdot5\cdot2$, $25\cdot2\cdot2\cdot2$, $\dots$, $200$) gives a number larger than $5\times10^7$. (Merging two factors $f,g$ into $fg$ always replaces $p^{g-1}$ by a factor at least $2^{f(g-1)}$, which here is much larger.) So the smallest number is $498960=2^4\cdot3^4\cdot5\cdot7\cdot11$, and its digit sum is $4+9+8+9+6+0=\boxed{36}$.
